If $\dfrac{a}{b}<\dfrac{c}{d}$ and $bc-ad=1$, any rational $\dfrac{p}{q}$ strictly between them has
\[
q=d(bp-aq)+b(cq-dp)\ge b+d,
\]
because the two parenthesized quantities are positive integers.

For the first step, denominators $1,2$ give no fraction strictly between $1$ and $\sqrt2$, while $\dfrac{4}{3}$ works. Thus $(p_1,q_1)=(4,3)$. Since $\sqrt2<\dfrac{3}{2}$, the neighbouring bounds $\dfrac{4}{3}$ and $\dfrac{3}{2}$ force $q_2\ge5$. Their mediant $\dfrac{7}{5}<\sqrt2$ works, giving $(p_2,q_2)=(7,5)$.

Next $\sqrt2<\dfrac{17}{12}$, and $17\cdot5-7\cdot12=1$. Hence every fraction between $\dfrac{7}{5}$ and $\sqrt2$ has denominator at least $5+12=17$. The mediant $\dfrac{24}{17}$ works, since $24^2=576<578=2\cdot17^2$. Therefore $\boxed{q_3=17}$.
